\chapter{The Triangle-Gluing Theorem and the Real MUB Obstruction}
\label{ch:real-mub-obstruction}

\section{Exact Closure Upstairs}

Represent each of the three contexts $Z,X,Y$ by a real angle: fix $a_Z = 0$ by
an overall choice of coordinates (always available, and immaterial to any
observable transition matrix), and write $a_X = \alpha$, $a_Y = \beta$ for
real numbers. With $L_{AB} := R(a_B - a_A)$ as in
Chapter~\ref{ch:context-graphs},
\begin{equation}
L_{ZX} L_{XY} L_{YZ} = R(\alpha) R(\beta-\alpha) R(-\beta) = R(0) = I
\label{eq:exact-closure-explicit}
\end{equation}
exactly, with no quotient or approximation: this is the additive law for
rotation angles, true for \emph{every} choice of $\alpha,\beta \in \R$.

\section{The Fibre of the Probability Projection}

For a labeled real angle $a$, the diagonal transition probability relative to
the fixed reference at angle $0$ is $t(a) = \cos^2 a$, and the canonical
unsigned observable is $\theta(a) := \min\big(\arccos\sqrt{t(a)},\ \pi/2 -
\arccos\sqrt{t(a)}\big) \in [0,\pi/4]$, invariant under the label-swap gauge
of \S\ref{sec:gauge-single-frame}. Restrict attention to $\alpha,\beta \in
[0,\pi/2]$, where $\arccos\sqrt{\cos^2\theta} = \theta$ identically, so
$\theta(\alpha)=\alpha$ before folding and $\theta_{ZX} := \theta(\alpha)$,
$\theta_{ZY} := \theta(\beta)$ may be taken as the raw angles themselves on
this range.

\begin{lemma}[Closed-form composition identity]
\label{lem:closed-form}
\index{Closed-form composition identity (lemma)}
For $\theta_1,\theta_2 \in [0,\pi/2]$,
\begin{equation}
1 - \cos^2\theta_1 - \cos^2\theta_2 + 2\cos^2\theta_1\cos^2\theta_2 \pm
2\sqrt{\cos^2\theta_1\cos^2\theta_2\sin^2\theta_1\sin^2\theta_2} \ =\
\cos^2(\theta_1 \mp \theta_2).
\end{equation}
\end{lemma}

\begin{proof}
Expand $\cos^2(\theta_1-\theta_2) = (\cos\theta_1\cos\theta_2 +
\sin\theta_1\sin\theta_2)^2 = \cos^2\theta_1\cos^2\theta_2 +
\sin^2\theta_1\sin^2\theta_2 + 2\sin\theta_1\cos\theta_1\sin\theta_2\cos\theta_2$.
Since $\theta_1,\theta_2 \in [0,\pi/2]$, $\sin\theta_i,\cos\theta_i \ge 0$, so
$\sin\theta_i\cos\theta_i = \sqrt{\cos^2\theta_i\sin^2\theta_i}$, and
$\sin^2\theta_1\sin^2\theta_2 = (1-\cos^2\theta_1)(1-\cos^2\theta_2)$; expanding
the latter and collecting terms gives the stated identity with the $+$ sign,
matching $\cos^2(\theta_1-\theta_2)$. The $-$ sign case follows identically
using $\cos^2(\theta_1+\theta_2) = (\cos\theta_1\cos\theta_2 -
\sin\theta_1\sin\theta_2)^2$.
\end{proof}

Writing $t_i = \cos^2\theta_i$, Lemma~\ref{lem:closed-form} says exactly that
$\cos^2(\beta-\alpha)$, computed from $t_{ZX}=\cos^2\alpha$ and
$t_{ZY}=\cos^2\beta$ alone, is not single-valued: it is one of the two roots
\begin{equation}
t_{XY} \in \big\{\, \cos^2(\theta_{ZX}-\theta_{ZY}),\ \cos^2(\theta_{ZX}+\theta_{ZY})
\,\big\}
\end{equation}
according to the (unobservable) relative sign of $\sin\alpha$ and $\sin\beta$.
This is the precise fibre statement in place of the earlier proof sketch: the
projection loses exactly one bit — a relative sign — beyond the diagonal
probabilities $t_{ZX}, t_{ZY}$ themselves.

\section{Descent: the Triangle-Gluing Theorem}

\begin{theorem}[Triangle gluing in $\R^2$]
\label{thm:triangle-gluing}
\index{Triangle gluing in R2 (theorem)}
Let $\theta_{ZX},\theta_{ZY},\theta_{XY} \in [0,\pi/4]$ be the canonical
unsigned transition angles of three pairwise $2\times2$ bistochastic matrices
on contexts $Z,X,Y$. A real assignment of frames to $Z,X,Y$ jointly realizing
all three exists if and only if
\begin{equation}
\theta_{XY} \equiv \pm\theta_{ZX} \pm \theta_{ZY} \pmod{\pi/2}
\label{eq:triangle-gluing-relation}
\end{equation}
for some choice of signs.
\end{theorem}

\begin{proof}
($\Rightarrow$, necessity.) Immediate from \eqref{eq:exact-closure-explicit}
and Lemma~\ref{lem:closed-form}: any real realization has $\theta_{XY}$ equal,
before its own canonical folding, to $\theta_{ZX}\pm\theta_{ZY}$ for one of the
two signs; folding to $[0,\pi/4]$ turns equality into the congruence
\eqref{eq:triangle-gluing-relation}.

($\Leftarrow$, sufficiency.) Suppose \eqref{eq:triangle-gluing-relation} holds
for some choice of signs. Fix $a_Z = 0$, $a_X = \theta_{ZX}$, and
\begin{equation}
a_Y = \begin{cases} \theta_{ZY} & \text{if the relation holds with } \theta_{XY}
\equiv \theta_{ZX} - \theta_{ZY} \pmod{\pi/2}, \\ -\theta_{ZY} & \text{if it
holds with } \theta_{XY} \equiv \theta_{ZX} + \theta_{ZY} \pmod{\pi/2}. \end{cases}
\end{equation}
By construction $\cos^2 a_X = \cos^2\theta_{ZX}$ and $\cos^2 a_Y =
\cos^2(\pm\theta_{ZY}) = \cos^2\theta_{ZY}$, reproducing $B_{ZX}$ and $B_{ZY}$
exactly; and $a_Y - a_X$ equals $-(\theta_{ZX}+\theta_{ZY})$ or
$\theta_{ZY}-\theta_{ZX}$ respectively, whose canonical fold is $\theta_{XY}$
by hypothesis, reproducing $B_{XY}$. This exhibits an explicit real frame
assignment realizing all three matrices.
\end{proof}

\begin{remark}
This theorem is the general statement; mutual unbiasedness is the sharpest
special case, not the only one. But "sharpest" should not be read as
"rare": for \emph{independently and continuously distributed} edge angles,
realizability itself is the exceptional condition. Fixing any two of the
three angles, Theorem~\ref{thm:triangle-gluing} shows the third is
constrained to one of at most two specific values out of the continuum
$[0,\pi/4]$ — a codimension-one condition. If all three angles are drawn
independently from a continuous distribution, the probability of hitting
that condition exactly is zero. Realizability is not generic; it is the
non-generic, measure-zero case, true for every graph containing a cycle,
not a special feature of the MUB point. What is special about mutual
unbiasedness is symmetry, not scarcity: it is the unique maximally
symmetric point at which the (already generic) failure occurs, not the
rare exception to an otherwise generic success.
\end{remark}

\section{The Mutually Unbiased Corollary}
\label{sec:mub-corollary}

\begin{corollary}[Real MUB obstruction, $d=2$]
\label{cor:real-mub}
\index{Real MUB obstruction, d=2 (corollary)}
There do not exist three pairwise mutually unbiased orthonormal bases of $\R^2$.
\end{corollary}

\begin{proof}
Set $\theta_{ZX} = \theta_{ZY} = \theta_{XY} = \pi/4$ (mutual unbiasedness).
Theorem~\ref{thm:triangle-gluing} requires $\pi/4 \equiv \pm\pi/4 \pm \pi/4
\pmod{\pi/2}$. The right-hand side takes only the values $0$ or $\pi/2$, both
$\equiv 0 \pmod{\pi/2}$, never $\pi/4$. Contradiction. (Equivalently, via
Lemma~\ref{lem:closed-form} directly: with $\theta_1=\theta_2=\pi/4$,
$\cos^2(\theta_1\mp\theta_2) \in \{\cos^2 0,\ \cos^2(\pi/2)\} = \{1,0\}$, never
$\cos^2(\pi/4) = 1/2$.)
\end{proof}

\begin{scopebox}
\scopetitle{}This rules out $\R^2$ for one specific configuration — three
mutually unbiased contexts — not for real-valued physics in general. It
does not show $\R$ fails whenever a cycle is present (most cycle data
\emph{is} real-realizable; \S\ref{sec:mub-corollary}'s remark below makes
this explicit), and it does not rule out an enlarged real state space or
a quaternionic alternative, both explicitly untouched here (this
chapter's own closing ledger says so). Treat "real numbers cannot do
this" as scoped to this exact triangle, this exact dimension, this exact
data — not as a general verdict on $\R$.
\end{scopebox}

\begin{remark}
Corollary~\ref{cor:real-mub} is the $d=2$ specialization of the general real-MUB
bound $N_{\R}(d) \le \lfloor d/2 \rfloor + 1$ of Boykin, Sitharam, Tarifi, and
Wocjan~\cite{boykin}. The bound is not new; what the reconstruction contributes
is the interpretation: every edge of the context triangle is individually
orthostochastically admissible, yet the triangle as a whole has no realization in
$\R^2$, while $\C^2$ realizes all three simultaneously via
\begin{equation}
Z = \left\{\binom{1}{0}, \binom{0}{1}\right\}, \quad
X = \tfrac{1}{\sqrt2}\left\{\binom{1}{1}, \binom{1}{-1}\right\}, \quad
Y = \tfrac{1}{\sqrt2}\left\{\binom{1}{i}, \binom{1}{-i}\right\}.
\end{equation}
\end{remark}

\begin{ledgerentry}[End of Part I, single-system results]
\textbf{Established:} edgewise real admissibility does not imply global real
admissibility (Theorem~\ref{thm:triangle-gluing}, proved in both directions:
necessity from the closed-form fibre identity of
Lemma~\ref{lem:closed-form}, sufficiency by explicit construction); the
mutually unbiased triple is the maximally inconsistent case
(Corollary~\ref{cor:real-mub}); this rules out $\R^2$ as a global lift for
three pairwise-MUB contexts without invoking composite systems.\\
\textbf{Not established:} that quantum mechanics must be complex in general;
enlarged real state spaces and quaternionic alternatives remain untouched by this
argument. The uniqueness question is deferred to Part II (composite systems,
local tomography).
\end{ledgerentry}

\begin{remark}[A connection noticed, not pursued]
The cycle-based obstruction of this chapter has the flavor of an early,
primitive version of holonomy — a closed loop carrying information no
single edge does, the same shape as Berry phase or a $\Z_2$-valued
cocycle failing to trivialize. This book does not develop that connection:
Part IV pursued higher-dimensional and quaternionic witnesses instead of a
Berry-phase chapter, so the analogy above should be read as exactly that
— an analogy noticed in passing, not a bundle/cocycle formulation worked
out rigorously anywhere in this book. It is recorded here as a genuine
open direction rather than silently dropped.
\end{remark}
